\documentclass[11pt,letterpaper]{article}

% Compile from this file's own folder (CHE374/problems/).
% Shared formatting for course notes. Loaded by each course's main.tex.
\usepackage[T1]{fontenc}
\usepackage[utf8]{inputenc}
\usepackage{lmodern}
\linespread{1.15}
\usepackage[margin=1in]{geometry}
\usepackage{amsmath,amssymb,amsthm}
\usepackage{mathtools}
\usepackage{graphicx}
\usepackage{booktabs}
\usepackage{enumitem}
\usepackage{xcolor}
\usepackage{hyperref}
\usepackage{tikz}
\usetikzlibrary{decorations.pathreplacing,positioning}
\usepackage{xparse}
\usepackage{etoolbox}

\hypersetup{colorlinks=true,linkcolor=blue!50!black,urlcolor=blue!50!black}
\setlength{\parindent}{0pt}
\setlength{\parskip}{0.5em}
\setlist{nosep}

% Unnumbered section headings: entries show as "Week 1", not "1 Week 1",
% both in the body and in the table of contents.
\setcounter{secnumdepth}{0}

% Each \section (i.e. each week's entry) starts on its own new page. The
% \newpage that used to follow \tableofcontents in each course's main.tex
% is redundant now and has been removed, so week 1 isn't pushed to a blank
% page by this.
\pretocmd{\section}{\newpage}{}{\errmessage{Failed to patch \string\section}}

% Tracks whether \section was the most recent sectioning/lecture command,
% so a \newlecture immediately following a \section (the usual "week's
% first lecture" case) doesn't also throw in its own \newpage, which would
% otherwise leave a blank page between the two. Set true by \section, and
% cleared by \newlecture whether or not it was true, so it only fires for
% a \newlecture that directly follows a \section (nothing else in between).
\newbool{newlecture@justsection}
\pretocmd{\section}{\booltrue{newlecture@justsection}}{}{\errmessage{Failed to patch \string\section}}

% Custom head format: bold, unparenthesized title, with a \quad gap after
% the number (e.g. "Theorem 3.2  Distribution" instead of "Theorem 3.2
% (Distribution)."), reusing amsthm's usual spacing/punctuation otherwise.
% When a title is given, the body text starts on its own new line below the
% head instead of running inline after it; untitled (auto-numbered-only)
% instances stay inline as usual. Body text is upright (not italic) for all
% six environments. Two variants share this head format but differ in
% end-of-head punctuation: theorem/lemma/proposition/exercise keep a
% trailing period, definition/example drop it (no period after the title).
\makeatletter
\newcommand{\thmpunct}{.}
\newtheoremstyle{boldhead}{}{}{\normalfont}{}{\bfseries}{\thmpunct}{5pt plus 1pt minus 1pt}%
    {\thmname{#1}\thmnumber{\@ifnotempty{#1}{ }#2}\thmnote{\@ifnotempty{#3}{\quad#3}}%
     \@ifnotempty{#3}{\def\thmheadnl{\newline}\thm@headsep\z@skip}{}}
\newtheoremstyle{boldheadnp}{}{}{\normalfont}{}{\bfseries}{}{5pt plus 1pt minus 1pt}%
    {\thmname{#1}\thmnumber{\@ifnotempty{#1}{ }#2}\thmnote{\@ifnotempty{#3}{\quad#3}}%
     \@ifnotempty{#3}{\def\thmheadnl{\newline}\thm@headsep\z@skip}{}}
\makeatother

\theoremstyle{boldhead}
\newtheorem{theorem}{Theorem}
\newtheorem{lemma}[theorem]{Lemma}
\newtheorem{proposition}[theorem]{Proposition}
\newtheorem{corollary}[theorem]{Corollary}
\newtheorem{exercise}[theorem]{Exercise}
\theoremstyle{boldheadnp}
\newtheorem{definition}[theorem]{Def.}
\newtheorem{example}[theorem]{Example}
\theoremstyle{remark}
\newtheorem*{remark}{Remark}

% \begin{theorem}[Title][2.7]: both optional arguments are new on top of
% the environments' usual single bracket. The first argument still behaves
% exactly like plain amsthm's optional title (existing notes using
% \begin{example}[Some Name] etc. are unaffected). The second argument, if
% given, overrides the displayed number (e.g. to match a textbook's own
% numbering) instead of the shared theorem counter's automatic value; the
% counter itself still advances normally, so later auto-numbered
% environments are unaffected too.
\newcommand{\addtitleandnumber}[1]{%
    \csletcs{#1orig}{#1}%
    \csletcs{end#1orig}{end#1}%
    \RenewDocumentEnvironment{#1}{o o}{%
        \IfValueT{##2}{\ifblank{##2}{}{\renewcommand{\thetheorem}{##2}}}%
        \IfValueTF{##1}{\ifblank{##1}{\csname#1orig\endcsname}{\csname#1orig\endcsname[##1]}}{\csname#1orig\endcsname}%
    }{\csname end#1orig\endcsname}%
}
\addtitleandnumber{theorem}
\addtitleandnumber{lemma}
\addtitleandnumber{proposition}
\addtitleandnumber{corollary}
\addtitleandnumber{definition}
\addtitleandnumber{example}
\addtitleandnumber{exercise}

% Common mathematical notation.
\newcommand{\R}{\mathbb{R}}
\newcommand{\N}{\mathbb{N}}
\newcommand{\Z}{\mathbb{Z}}
\newcommand{\E}{\mathbb{E}}
\newcommand{\Prob}{\mathbb{P}}
\DeclareMathOperator{\Var}{Var}
\DeclareMathOperator{\Cov}{Cov}

% \newlecture[title][date]: bold, centered marker for the start of a new
% lecture. Both arguments are optional. Numbering runs continuously across
% the whole document (i.e. across all weeks' entries, not reset per
% \section), so the Nth \newlecture in the course is always "Lecture N".
% Each lecture starts on its own new page, except when it directly follows
% a \section heading (which already started a new page of its own), so the
% two stay together instead of leaving a blank page in between.
\newcounter{lecture}
\renewcommand{\thelecture}{\arabic{lecture}}
\NewDocumentCommand{\newlecture}{o o}{%
    \stepcounter{lecture}%
    \ifbool{newlecture@justsection}{}{\newpage}%
    \boolfalse{newlecture@justsection}%
    \par\vspace{1em}
    \begin{center}
        \textbf{\Large Lecture \thelecture\IfValueT{#1}{: #1}}\\[2pt]
        \IfValueT{#2}{\normalsize #2}
    \end{center}
    \vspace{0.5em}%
}

% Handwritten-note-style helpers: red definition labels, blue margin
% annotations, yellow highlights, underlined mini-headers.
\newcommand{\deflabel}[2]{\par\noindent{\color{red!70!black}\textbf{Def #1}}\quad\textbf{#2}\par\nobreak}
\newcommand{\exlabel}[2]{\par\noindent\textbf{Example #1}\quad\textbf{#2}\par\nobreak}
\newcommand{\heading}[1]{\par\noindent{\color{red!70!black}\textbf{#1}}\par}
\newcommand{\header}[1]{\par\vspace{0.5em}\noindent\textbf{#1}\par\vspace{-0.2em}}
\newcommand{\aside}[1]{\ifmmode{\color{blue!55!black}\text{\itshape #1}}\else{\color{blue!55!black}\itshape #1}\fi}
\newcommand{\hly}[1]{\colorbox{yellow!55}{$#1$}}
\newcommand{\hlt}[1]{\colorbox{yellow!55}{#1}}
% \boldmath makes any math the argument switches into (e.g. \rb{$r$}) use
% the bold math version automatically, so callers don't need to manually
% wrap math content in \boldsymbol{}.
\newcommand{\rb}[1]{{\color{red!70!black}\boldmath\textbf{#1}}}
\renewcommand{\b}[1]{\textbf{#1}}
\definecolor{navy}{RGB}{0,0,128}
\newcommand{\dbb}[1]{{\color{navy}\textbf{#1}}}
\newcommand{\db}[1]{{\color{navy}#1}}


\title{CHE374 --- Problem Set 1 Solutions}
\author{Bonnie Zhuang}
\date{Fall 2026}

\begin{document}
\maketitle

\section{Problem Set 1}

\subsection{Question 1}
\begin{quote}
    \itshape
    A new steel stirred tank reactor costs \$8,000, plus \$1,200 for
    installation. The company has an old stainless steel unit in the
    warehouse that could be used instead. The old reactor was used for
    4 years but will last the same length of time as the new one; it
    cost \$16,000 when purchased. Installing the old reactor would cost
    \$2,500. The old reactor has a used market value of \$6,800.
    \\ a) Determine the best alternative.
    \\ b) What market value would make you indifferent between the two?
\end{quote}

\textbf{Solution:} The \$16,000 original purchase price of the old
reactor is a \textbf{sunk cost} --- already spent, and irrelevant to
the decision going forward. What matters is the \emph{opportunity
cost} of using the old reactor instead of selling it at its market
value. Since both reactors have the same remaining life, we can
compare total up-front cost directly.
\begin{itemize}
    \item New reactor: $8{,}000 + 1{,}200 = \$9{,}000$.
    \item Old reactor: $6{,}800 + 2{,}500 = \$9{,}300$ (market value
    forgone, plus installation).
\end{itemize}
a) Since $\$9{,}000 < \$9{,}300$, buy the \textbf{new reactor}.

b) Let $X$ be the old reactor's market value. Indifference means
\[X + 2{,}500 = 9{,}000 \implies X = \$6{,}500.\]
Since the actual market value (\$6,800) is above this breakeven point,
the old reactor is relatively too valuable to give up for free ---
confirming the new reactor is the better buy.

\subsection{Question 2}
\begin{quote}
    \itshape
    A restaurant is upgrading its cash registers. Company X quotes
    \$7{,}500 up front, saving \$1{,}700/year. Company Y quotes
    \$9{,}000 up front, saving \$2{,}200/year.
    \\ a) How long to recover the initial investment for each?
    \\ b) List at least one issue with this type of comparison.
\end{quote}

\textbf{Solution:}
\\ a) Payback period $=$ initial cost / annual savings.
\[\text{Company X: } \frac{7{,}500}{1{,}700} \approx 4.41 \text{ years}, \qquad \text{Company Y: } \frac{9{,}000}{2{,}200} \approx 4.09 \text{ years}.\]
Company Y pays back faster despite the larger upfront cost.

b) Simple payback analysis has (at least) two issues:
\begin{itemize}
    \item It ignores the \textbf{time value of money} --- future
    savings are not discounted back to present worth.
    \item It ignores everything that happens \emph{after} the payback
    period (e.g.\ total savings over the equipment's full life, or
    differences in equipment lifespan/reliability between options).
\end{itemize}

\subsection{Question 3}
\begin{quote}
    \itshape
    Find $\displaystyle\lim_{n\to\infty}\left(1+\frac{r}{n}\right)^n$.
\end{quote}

\textbf{Solution:} This is exactly the continuous-compounding limit
from the notes (with the ``$-1$'' dropped):
\[\lim_{n\to\infty}\left(1+\frac{r}{n}\right)^n = e^r.\]

\subsection{Question 4}
\begin{quote}
    \itshape
    An 8-year \$100,000 loan charges an effective 3.4\% interest per
    year. Find the nominal annual rate for: a) semi-annual, b) monthly,
    c) daily, d) continuous compounding. e) How much is owed after 4
    years in each case?
\end{quote}

\textbf{Solution:} Solve $i_e = \left(1+\frac{r}{m}\right)^m - 1 = 0.034$
for $r$, i.e.\ $r = m\left[(1.034)^{1/m}-1\right]$.
\begin{itemize}
    \item a) $m=2$: $r = 2\left[(1.034)^{1/2}-1\right] \approx 3.372\%$.
    \item b) $m=12$: $r = 12\left[(1.034)^{1/12}-1\right] \approx 3.348\%$.
    \item c) $m=365$: $r = 365\left[(1.034)^{1/365}-1\right] \approx 3.343\%$.
    \item d) Continuous: $r = \ln(1.034) \approx 3.343\%$ (matches
    $e^r - 1 = i_e$ from the notes).
\end{itemize}

e) Since every case was solved to give the \emph{same} effective
annual rate (3.4\%), the amount owed after 4 years is identical
regardless of compounding frequency:
\[F_4 = 100{,}000(1.034)^4 \approx \$114{,}309.46 \quad \text{in all four cases.}\]
This is the point of the effective rate: it's the true, comparable
growth rate, independent of how compounding is described.

\subsection{Question 5}
\begin{quote}
    \itshape
    A \$100 T-bill (face value paid at maturity, no coupons) sells at a
    continuously compounding yield of 2.3\%/year.
    \\ a) Effective annual yield with yearly compounding?
    \\ b) Price if it matures in 6 months? c) In 9 months?
    \\ d) If the yield drops 30 bps (to 2\%), how do the prices in b)
    and c) change?
\end{quote}

\textbf{Solution:} For continuous compounding, price $= Fe^{-rt}$
where $F=\$100$ is the face value and $t$ is time to maturity in years.

a) $i_e = e^{0.023} - 1 \approx 2.327\%$.

b) $t = 0.5$: $\text{Price} = 100\,e^{-0.023(0.5)} \approx \$98.86$.

c) $t = 0.75$: $\text{Price} = 100\,e^{-0.023(0.75)} \approx \$98.29$.

d) With $r = 2\%$:
\[\text{6-month: } 100\,e^{-0.02(0.5)} \approx \$99.00 \quad (+\$0.15), \qquad \text{9-month: } 100\,e^{-0.02(0.75)} \approx \$98.51 \quad (+\$0.22).\]
Both prices \emph{increase} when the yield drops --- price and yield
move in opposite directions for a zero-coupon bond --- and the
9-month (longer maturity) price is more sensitive to the yield
change than the 6-month price.

\subsection{Question 6}
\begin{quote}
    \itshape
    Interest rate conversion table (assume 252 trading days/year where
    a daily rate is involved).
\end{quote}

\textbf{Solution:}
\begin{enumerate}
    \item \textbf{Given} 8\% effective/year $\to$ \textbf{monthly rate,
    monthly compounding}: $i_s = (1.08)^{1/12} - 1 \approx 0.643\%$/month.
    \item \textbf{Given} 3.5\%/year nominal, daily compounding ($m=252$)
    $\to$ \textbf{effective rate/year}: $i_e = \left(1+\frac{0.035}{252}\right)^{252}-1 \approx 3.562\%$.
    \item \textbf{Given} 4\%/year nominal, quarterly compounding
    ($m=4$) $\to$ \textbf{continuous rate/year}: $r_c = 4\ln(1.01) \approx 3.980\%$.
    \item \textbf{Given} 1.5\%/month, monthly compounding $\to$
    \textbf{continuous rate/quarter}: 3 months per quarter, so
    $r_c = 3\ln(1.015) \approx 4.467\%$/quarter.
    \item \textbf{Given} 1.2\%/quarter, monthly compounding ($m=3$
    months/quarter, so $0.4\%$/month) $\to$ \textbf{interest per 3
    years, semi-annual compounding}: the equivalent effective rate is
    $(1.004)^{36}-1 \approx 15.455\%$ total over 3 years (equivalently,
    an effective semi-annual rate of $(1.004)^6-1 \approx 2.423\%$,
    compounded over the 6 half-years in 3 years).
\end{enumerate}

\end{document}
